\documentclass[11pt]{article}
% ---------------------------------------------------------------
% R4.02 Qualite de developpement -- preambule commun des sujets de TP et des corrections
% Article clair, fond blanc, accents de couleur coherents avec les slides
% Charte reprise de R3.07 (meme palette IUT que les diapositives)
% Utilisation :
%   \documentclass[11pt]{article}
%   % ---------------------------------------------------------------
% R4.02 Qualite de developpement -- preambule commun des sujets de TP et des corrections
% Article clair, fond blanc, accents de couleur coherents avec les slides
% Charte reprise de R3.07 (meme palette IUT que les diapositives)
% Utilisation :
%   \documentclass[11pt]{article}
%   % ---------------------------------------------------------------
% R4.02 Qualite de developpement -- preambule commun des sujets de TP et des corrections
% Article clair, fond blanc, accents de couleur coherents avec les slides
% Charte reprise de R3.07 (meme palette IUT que les diapositives)
% Utilisation :
%   \documentclass[11pt]{article}
%   \input{../../latex/common/r402-tp-preamble.tex}
%   \begin{document}
%   \rtpheader{1}{Titre du TP}
% ---------------------------------------------------------------
\usepackage[T1]{fontenc}
\usepackage[utf8]{inputenc}
\usepackage[french]{babel}
\usepackage[scale=0.92]{tgheros}
\usepackage{tgpagella}
\usepackage[a4paper, margin=2.2cm, headheight=14pt]{geometry}
\usepackage{graphicx}
\usepackage{xcolor}
\usepackage{listings}
\usepackage[skins,listings,breakable]{tcolorbox}
\usepackage{enumitem}
\usepackage{booktabs}
\usepackage{array}
\usepackage{multirow}
\usepackage{tikz}
\usetikzlibrary{shapes, shapes.geometric, arrows, arrows.meta, positioning, calc}
\usepackage[colorlinks=true, linkcolor=qaaccentdark, urlcolor=qaaccentdark]{hyperref}
\usepackage{fancyhdr}
\usepackage{lastpage}
\usepackage{titlesec}
\usepackage{etoolbox}
\usepackage{eurosym}

% ---------------------------------------------------------------
% MENTION D'ELABORATION (discrete : gris clair, petit corps)
% ---------------------------------------------------------------
\definecolor{qaaimention}{HTML}{A3AAB5}
\newcommand{\raimention}{{\scriptsize\color{qaaimention}Supports \'elabor\'es avec l'aide de Claude Code Opus 5.5}}

% ---------------------------------------------------------------
% PALETTE -- version imprimable des couleurs des diapositives
% ---------------------------------------------------------------
\definecolor{qaaccentdark}{HTML}{7A1680}   % violet du logo IUT, fonce pour rester lisible sur blanc
\definecolor{qaaccent2dark}{HTML}{9B4B94}  % violet clair du logo IUT (accent secondaire)
\definecolor{qapanellight}{HTML}{F4F6F8}   % fond de bloc tres clair
\definecolor{qalinelight}{HTML}{D8DEE4}    % filets
\definecolor{qatextdark}{HTML}{1A1D23}     % texte principal
\definecolor{qamuteddark}{HTML}{5B6472}    % texte secondaire
\definecolor{qacodekwl}{HTML}{7A1680}      % code : mots-cles
\definecolor{qacodestrl}{HTML}{2E7D32}     % code : chaines
\definecolor{qacodecoml}{HTML}{6B7280}     % code : commentaires
\definecolor{qacodeannl}{HTML}{B26A00}     % code : annotations / javadoc

\color{qatextdark}

\setlength{\parskip}{4pt}
\setlength{\parindent}{0pt}

% ---------------------------------------------------------------
% TITRES DE SECTION -- texte blanc sur bande de couleur pleine largeur
% (gere aussi les variantes etoilees \section*)
% ---------------------------------------------------------------
\newcommand{\qasectionband}[2]{%
  \noindent\colorbox{#1}{%
    \parbox[c]{\dimexpr\linewidth-2\fboxsep\relax}{%
      \color{white}#2%
    }%
  }%
}
\newcommand{\secformat}[1]{\qasectionband{qaaccentdark}{\bfseries\Large\thesection\quad #1}}
\newcommand{\secformatnn}[1]{\qasectionband{qaaccentdark}{\bfseries\Large #1}}
\newcommand{\subsecformat}[1]{\qasectionband{qaaccent2dark}{\bfseries\large\thesubsection\quad #1}}
\newcommand{\subsecformatnn}[1]{\qasectionband{qaaccent2dark}{\bfseries\large #1}}

\titleformat{\section}{\normalfont}{}{0pt}{\secformat}
\titleformat{name=\section,numberless}{\normalfont}{}{0pt}{\secformatnn}
\titleformat{\subsection}{\normalfont}{}{0pt}{\subsecformat}
\titleformat{name=\subsection,numberless}{\normalfont}{}{0pt}{\subsecformatnn}
\titleformat{\subsubsection}{\normalfont\bfseries\color{qaaccentdark}}{\thesubsubsection}{0.6em}{}
\titlespacing*{\section}{0pt}{14pt}{8pt}
\titlespacing*{\subsection}{0pt}{10pt}{6pt}
\titlespacing*{\subsubsection}{0pt}{8pt}{4pt}

% ---------------------------------------------------------------
% LISTINGS -- theme clair (lstlisting reste utilisable tel quel)
% \lstcode : police + couleur de base, a reutiliser dans les surcharges
% ---------------------------------------------------------------
\newcommand{\lstcode}{\ttfamily\color{qatextdark}}

\lstdefinelanguage{TypeScript}{
  keywords={typeof, new, true, false, catch, function, return, null, switch, var,
    if, in, while, do, else, case, break, const, let, interface, class, extends,
    implements, enum, export, import, this, super, constructor, public, private,
    protected, static, async, await, throw, throws, try, finally, void},
  ndkeywords={boolean, any, number, string, int, long, double},
  sensitive=true,
  comment=[l]{//},
  morecomment=[s]{/*}{*/},
  morestring=[b]',
  morestring=[b]"
}

\lstdefinestyle{r402code}{
  basicstyle=\lstcode\footnotesize,
  keywordstyle=\color{qacodekwl}\bfseries,
  ndkeywordstyle=\color{qaaccent2dark}\bfseries,
  stringstyle=\color{qacodestrl},
  commentstyle=\color{qacodecoml}\itshape,
  morecomment=[s][\color{qacodeannl}]{/**}{*/},
  backgroundcolor=\color{qapanellight},
  rulecolor=\color{qalinelight},
  frame=single,
  framesep=4pt,
  xleftmargin=6pt,
  xrightmargin=6pt,
  numbers=left,
  numberstyle=\tiny\color{qamuteddark},
  numbersep=8pt,
  keepspaces=true,
  columns=fullflexible,
  showstringspaces=false,
  showspaces=false,
  showtabs=false,
  breaklines=true,
  breakatwhitespace=false,
  captionpos=b,
  tabsize=2,
  escapeinside={(*@}{@*)},
  % accents UTF-8 dans le code (listings ne gere pas l'UTF-8 multi-octets)
  literate={é}{{\'e}}1 {è}{{\`e}}1 {ê}{{\^e}}1 {à}{{\`a}}1 {â}{{\^a}}1
           {ù}{{\`u}}1 {û}{{\^u}}1 {ô}{{\^o}}1 {î}{{\^i}}1 {ç}{{\c c}}1
           {É}{{\'E}}1 {À}{{\`A}}1 {€}{{\euro{}}}1 {…}{{\ldots}}1,
}
\lstset{style=r402code}

% ---------------------------------------------------------------
% BOITES : exercice, code, information
% ---------------------------------------------------------------
\tcbset{
  exobox/.style={colback=qapanellight, colframe=qaaccentdark, boxrule=0.6pt,
    arc=2pt, left=8pt, right=8pt, top=6pt, bottom=6pt,
    coltitle=white, colbacktitle=qaaccentdark, fonttitle=\bfseries,
    boxsep=2pt, breakable},
  infobox/.style={colback=white, colframe=qaaccent2dark, boxrule=0.6pt,
    leftrule=3pt, arc=0pt, left=8pt, right=8pt, top=5pt, bottom=5pt,
    coltitle=qaaccent2dark, colbacktitle=white, fonttitle=\bfseries,
    breakable},
  codebox/.style={colback=qatextdark, colframe=qalinelight, boxrule=0.4pt,
    arc=2pt, left=8pt, right=8pt, top=5pt, bottom=5pt, listing only,
    listing options={basicstyle=\ttfamily\small\color{white},
      breaklines=true, showstringspaces=false, columns=fullflexible,
      keepspaces=true, frame=none, numbers=none, backgroundcolor={}}, breakable}
}
\newtcolorbox[auto counter]{exobox}[2][]{exobox,
  title={Exercice~\thetcbcounter\ifstrempty{#2}{}{ --- #2}}, #1}
\newtcolorbox{infobox}[1][\`A retenir]{infobox, title={#1}}
\newtcblisting{codebox}{codebox}

% ---------------------------------------------------------------
% EN-TETE DE PAGE (fil d'Ariane + pagination)
% ---------------------------------------------------------------
\newcommand{\qaheadtitlevalue}{}
\newcommand{\qaheadtitle}[1]{\renewcommand{\qaheadtitlevalue}{#1}}

\pagestyle{fancy}
\fancyhf{}
\renewcommand{\headrulewidth}{0.4pt}
\fancyhead[L]{\small\color{qamuteddark}R4.02 --- \qaheadtitlevalue}
\fancyhead[R]{\small\color{qamuteddark}\thepage\,/\,\pageref{LastPage}}
\fancyfoot[C]{\footnotesize\color{qamuteddark}IUT de Vannes --- Universit\'e Bretagne Sud}
\fancypagestyle{plain}{\fancyhf{}\renewcommand{\headrulewidth}{0pt}%
  \fancyfoot[C]{\footnotesize\color{qamuteddark}IUT de Vannes --- Universit\'e Bretagne Sud}}

% ---------------------------------------------------------------
% BANDEAU DE TITRE (pas une page dediee, pour rester compact)
% \rtpheader{numero ou intitule court}{titre du TP}
% ---------------------------------------------------------------
\newcommand{\ryear}{2026/2027}
\newcommand{\rtpheader}[2]{%
  \thispagestyle{plain}%
  \begin{minipage}[t]{0.68\textwidth}
    {\small\color{qamuteddark}
      Universit\'e Bretagne Sud\\
      IUT de Vannes\\
      D\'epartement Informatique --- \ryear}
  \end{minipage}%
  \hfill
  \begin{minipage}[t]{0.28\textwidth}
    \raggedleft
    \includegraphics[width=2.6cm]{iut_logo.jpg}
  \end{minipage}

  \vspace{6mm}
  \begin{center}
    {\large\bfseries\color{qatextdark} BUT 2 --- R4.02. Qualit\'e de d\'eveloppement}\\[2mm]
    {\normalsize\color{qamuteddark} TP #1}\\[1mm]
    {\Large\bfseries\color{qaaccentdark} #2}\\[2mm]
    \raimention
  \end{center}
  \vspace{4mm}
  \qaheadtitle{#2}
}

%   \begin{document}
%   \rtpheader{1}{Titre du TP}
% ---------------------------------------------------------------
\usepackage[T1]{fontenc}
\usepackage[utf8]{inputenc}
\usepackage[french]{babel}
\usepackage[scale=0.92]{tgheros}
\usepackage{tgpagella}
\usepackage[a4paper, margin=2.2cm, headheight=14pt]{geometry}
\usepackage{graphicx}
\usepackage{xcolor}
\usepackage{listings}
\usepackage[skins,listings,breakable]{tcolorbox}
\usepackage{enumitem}
\usepackage{booktabs}
\usepackage{array}
\usepackage{multirow}
\usepackage{tikz}
\usetikzlibrary{shapes, shapes.geometric, arrows, arrows.meta, positioning, calc}
\usepackage[colorlinks=true, linkcolor=qaaccentdark, urlcolor=qaaccentdark]{hyperref}
\usepackage{fancyhdr}
\usepackage{lastpage}
\usepackage{titlesec}
\usepackage{etoolbox}
\usepackage{eurosym}

% ---------------------------------------------------------------
% MENTION D'ELABORATION (discrete : gris clair, petit corps)
% ---------------------------------------------------------------
\definecolor{qaaimention}{HTML}{A3AAB5}
\newcommand{\raimention}{{\scriptsize\color{qaaimention}Supports \'elabor\'es avec l'aide de Claude Code Opus 5.5}}

% ---------------------------------------------------------------
% PALETTE -- version imprimable des couleurs des diapositives
% ---------------------------------------------------------------
\definecolor{qaaccentdark}{HTML}{7A1680}   % violet du logo IUT, fonce pour rester lisible sur blanc
\definecolor{qaaccent2dark}{HTML}{9B4B94}  % violet clair du logo IUT (accent secondaire)
\definecolor{qapanellight}{HTML}{F4F6F8}   % fond de bloc tres clair
\definecolor{qalinelight}{HTML}{D8DEE4}    % filets
\definecolor{qatextdark}{HTML}{1A1D23}     % texte principal
\definecolor{qamuteddark}{HTML}{5B6472}    % texte secondaire
\definecolor{qacodekwl}{HTML}{7A1680}      % code : mots-cles
\definecolor{qacodestrl}{HTML}{2E7D32}     % code : chaines
\definecolor{qacodecoml}{HTML}{6B7280}     % code : commentaires
\definecolor{qacodeannl}{HTML}{B26A00}     % code : annotations / javadoc

\color{qatextdark}

\setlength{\parskip}{4pt}
\setlength{\parindent}{0pt}

% ---------------------------------------------------------------
% TITRES DE SECTION -- texte blanc sur bande de couleur pleine largeur
% (gere aussi les variantes etoilees \section*)
% ---------------------------------------------------------------
\newcommand{\qasectionband}[2]{%
  \noindent\colorbox{#1}{%
    \parbox[c]{\dimexpr\linewidth-2\fboxsep\relax}{%
      \color{white}#2%
    }%
  }%
}
\newcommand{\secformat}[1]{\qasectionband{qaaccentdark}{\bfseries\Large\thesection\quad #1}}
\newcommand{\secformatnn}[1]{\qasectionband{qaaccentdark}{\bfseries\Large #1}}
\newcommand{\subsecformat}[1]{\qasectionband{qaaccent2dark}{\bfseries\large\thesubsection\quad #1}}
\newcommand{\subsecformatnn}[1]{\qasectionband{qaaccent2dark}{\bfseries\large #1}}

\titleformat{\section}{\normalfont}{}{0pt}{\secformat}
\titleformat{name=\section,numberless}{\normalfont}{}{0pt}{\secformatnn}
\titleformat{\subsection}{\normalfont}{}{0pt}{\subsecformat}
\titleformat{name=\subsection,numberless}{\normalfont}{}{0pt}{\subsecformatnn}
\titleformat{\subsubsection}{\normalfont\bfseries\color{qaaccentdark}}{\thesubsubsection}{0.6em}{}
\titlespacing*{\section}{0pt}{14pt}{8pt}
\titlespacing*{\subsection}{0pt}{10pt}{6pt}
\titlespacing*{\subsubsection}{0pt}{8pt}{4pt}

% ---------------------------------------------------------------
% LISTINGS -- theme clair (lstlisting reste utilisable tel quel)
% \lstcode : police + couleur de base, a reutiliser dans les surcharges
% ---------------------------------------------------------------
\newcommand{\lstcode}{\ttfamily\color{qatextdark}}

\lstdefinelanguage{TypeScript}{
  keywords={typeof, new, true, false, catch, function, return, null, switch, var,
    if, in, while, do, else, case, break, const, let, interface, class, extends,
    implements, enum, export, import, this, super, constructor, public, private,
    protected, static, async, await, throw, throws, try, finally, void},
  ndkeywords={boolean, any, number, string, int, long, double},
  sensitive=true,
  comment=[l]{//},
  morecomment=[s]{/*}{*/},
  morestring=[b]',
  morestring=[b]"
}

\lstdefinestyle{r402code}{
  basicstyle=\lstcode\footnotesize,
  keywordstyle=\color{qacodekwl}\bfseries,
  ndkeywordstyle=\color{qaaccent2dark}\bfseries,
  stringstyle=\color{qacodestrl},
  commentstyle=\color{qacodecoml}\itshape,
  morecomment=[s][\color{qacodeannl}]{/**}{*/},
  backgroundcolor=\color{qapanellight},
  rulecolor=\color{qalinelight},
  frame=single,
  framesep=4pt,
  xleftmargin=6pt,
  xrightmargin=6pt,
  numbers=left,
  numberstyle=\tiny\color{qamuteddark},
  numbersep=8pt,
  keepspaces=true,
  columns=fullflexible,
  showstringspaces=false,
  showspaces=false,
  showtabs=false,
  breaklines=true,
  breakatwhitespace=false,
  captionpos=b,
  tabsize=2,
  escapeinside={(*@}{@*)},
  % accents UTF-8 dans le code (listings ne gere pas l'UTF-8 multi-octets)
  literate={é}{{\'e}}1 {è}{{\`e}}1 {ê}{{\^e}}1 {à}{{\`a}}1 {â}{{\^a}}1
           {ù}{{\`u}}1 {û}{{\^u}}1 {ô}{{\^o}}1 {î}{{\^i}}1 {ç}{{\c c}}1
           {É}{{\'E}}1 {À}{{\`A}}1 {€}{{\euro{}}}1 {…}{{\ldots}}1,
}
\lstset{style=r402code}

% ---------------------------------------------------------------
% BOITES : exercice, code, information
% ---------------------------------------------------------------
\tcbset{
  exobox/.style={colback=qapanellight, colframe=qaaccentdark, boxrule=0.6pt,
    arc=2pt, left=8pt, right=8pt, top=6pt, bottom=6pt,
    coltitle=white, colbacktitle=qaaccentdark, fonttitle=\bfseries,
    boxsep=2pt, breakable},
  infobox/.style={colback=white, colframe=qaaccent2dark, boxrule=0.6pt,
    leftrule=3pt, arc=0pt, left=8pt, right=8pt, top=5pt, bottom=5pt,
    coltitle=qaaccent2dark, colbacktitle=white, fonttitle=\bfseries,
    breakable},
  codebox/.style={colback=qatextdark, colframe=qalinelight, boxrule=0.4pt,
    arc=2pt, left=8pt, right=8pt, top=5pt, bottom=5pt, listing only,
    listing options={basicstyle=\ttfamily\small\color{white},
      breaklines=true, showstringspaces=false, columns=fullflexible,
      keepspaces=true, frame=none, numbers=none, backgroundcolor={}}, breakable}
}
\newtcolorbox[auto counter]{exobox}[2][]{exobox,
  title={Exercice~\thetcbcounter\ifstrempty{#2}{}{ --- #2}}, #1}
\newtcolorbox{infobox}[1][\`A retenir]{infobox, title={#1}}
\newtcblisting{codebox}{codebox}

% ---------------------------------------------------------------
% EN-TETE DE PAGE (fil d'Ariane + pagination)
% ---------------------------------------------------------------
\newcommand{\qaheadtitlevalue}{}
\newcommand{\qaheadtitle}[1]{\renewcommand{\qaheadtitlevalue}{#1}}

\pagestyle{fancy}
\fancyhf{}
\renewcommand{\headrulewidth}{0.4pt}
\fancyhead[L]{\small\color{qamuteddark}R4.02 --- \qaheadtitlevalue}
\fancyhead[R]{\small\color{qamuteddark}\thepage\,/\,\pageref{LastPage}}
\fancyfoot[C]{\footnotesize\color{qamuteddark}IUT de Vannes --- Universit\'e Bretagne Sud}
\fancypagestyle{plain}{\fancyhf{}\renewcommand{\headrulewidth}{0pt}%
  \fancyfoot[C]{\footnotesize\color{qamuteddark}IUT de Vannes --- Universit\'e Bretagne Sud}}

% ---------------------------------------------------------------
% BANDEAU DE TITRE (pas une page dediee, pour rester compact)
% \rtpheader{numero ou intitule court}{titre du TP}
% ---------------------------------------------------------------
\newcommand{\ryear}{2026/2027}
\newcommand{\rtpheader}[2]{%
  \thispagestyle{plain}%
  \begin{minipage}[t]{0.68\textwidth}
    {\small\color{qamuteddark}
      Universit\'e Bretagne Sud\\
      IUT de Vannes\\
      D\'epartement Informatique --- \ryear}
  \end{minipage}%
  \hfill
  \begin{minipage}[t]{0.28\textwidth}
    \raggedleft
    \includegraphics[width=2.6cm]{iut_logo.jpg}
  \end{minipage}

  \vspace{6mm}
  \begin{center}
    {\large\bfseries\color{qatextdark} BUT 2 --- R4.02. Qualit\'e de d\'eveloppement}\\[2mm]
    {\normalsize\color{qamuteddark} TP #1}\\[1mm]
    {\Large\bfseries\color{qaaccentdark} #2}\\[2mm]
    \raimention
  \end{center}
  \vspace{4mm}
  \qaheadtitle{#2}
}

%   \begin{document}
%   \rtpheader{1}{Titre du TP}
% ---------------------------------------------------------------
\usepackage[T1]{fontenc}
\usepackage[utf8]{inputenc}
\usepackage[french]{babel}
\usepackage[scale=0.92]{tgheros}
\usepackage{tgpagella}
\usepackage[a4paper, margin=2.2cm, headheight=14pt]{geometry}
\usepackage{graphicx}
\usepackage{xcolor}
\usepackage{listings}
\usepackage[skins,listings,breakable]{tcolorbox}
\usepackage{enumitem}
\usepackage{booktabs}
\usepackage{array}
\usepackage{multirow}
\usepackage{tikz}
\usetikzlibrary{shapes, shapes.geometric, arrows, arrows.meta, positioning, calc}
\usepackage[colorlinks=true, linkcolor=qaaccentdark, urlcolor=qaaccentdark]{hyperref}
\usepackage{fancyhdr}
\usepackage{lastpage}
\usepackage{titlesec}
\usepackage{etoolbox}
\usepackage{eurosym}

% ---------------------------------------------------------------
% MENTION D'ELABORATION (discrete : gris clair, petit corps)
% ---------------------------------------------------------------
\definecolor{qaaimention}{HTML}{A3AAB5}
\newcommand{\raimention}{{\scriptsize\color{qaaimention}Supports \'elabor\'es avec l'aide de Claude Code Opus 5.5}}

% ---------------------------------------------------------------
% PALETTE -- version imprimable des couleurs des diapositives
% ---------------------------------------------------------------
\definecolor{qaaccentdark}{HTML}{7A1680}   % violet du logo IUT, fonce pour rester lisible sur blanc
\definecolor{qaaccent2dark}{HTML}{9B4B94}  % violet clair du logo IUT (accent secondaire)
\definecolor{qapanellight}{HTML}{F4F6F8}   % fond de bloc tres clair
\definecolor{qalinelight}{HTML}{D8DEE4}    % filets
\definecolor{qatextdark}{HTML}{1A1D23}     % texte principal
\definecolor{qamuteddark}{HTML}{5B6472}    % texte secondaire
\definecolor{qacodekwl}{HTML}{7A1680}      % code : mots-cles
\definecolor{qacodestrl}{HTML}{2E7D32}     % code : chaines
\definecolor{qacodecoml}{HTML}{6B7280}     % code : commentaires
\definecolor{qacodeannl}{HTML}{B26A00}     % code : annotations / javadoc

\color{qatextdark}

\setlength{\parskip}{4pt}
\setlength{\parindent}{0pt}

% ---------------------------------------------------------------
% TITRES DE SECTION -- texte blanc sur bande de couleur pleine largeur
% (gere aussi les variantes etoilees \section*)
% ---------------------------------------------------------------
\newcommand{\qasectionband}[2]{%
  \noindent\colorbox{#1}{%
    \parbox[c]{\dimexpr\linewidth-2\fboxsep\relax}{%
      \color{white}#2%
    }%
  }%
}
\newcommand{\secformat}[1]{\qasectionband{qaaccentdark}{\bfseries\Large\thesection\quad #1}}
\newcommand{\secformatnn}[1]{\qasectionband{qaaccentdark}{\bfseries\Large #1}}
\newcommand{\subsecformat}[1]{\qasectionband{qaaccent2dark}{\bfseries\large\thesubsection\quad #1}}
\newcommand{\subsecformatnn}[1]{\qasectionband{qaaccent2dark}{\bfseries\large #1}}

\titleformat{\section}{\normalfont}{}{0pt}{\secformat}
\titleformat{name=\section,numberless}{\normalfont}{}{0pt}{\secformatnn}
\titleformat{\subsection}{\normalfont}{}{0pt}{\subsecformat}
\titleformat{name=\subsection,numberless}{\normalfont}{}{0pt}{\subsecformatnn}
\titleformat{\subsubsection}{\normalfont\bfseries\color{qaaccentdark}}{\thesubsubsection}{0.6em}{}
\titlespacing*{\section}{0pt}{14pt}{8pt}
\titlespacing*{\subsection}{0pt}{10pt}{6pt}
\titlespacing*{\subsubsection}{0pt}{8pt}{4pt}

% ---------------------------------------------------------------
% LISTINGS -- theme clair (lstlisting reste utilisable tel quel)
% \lstcode : police + couleur de base, a reutiliser dans les surcharges
% ---------------------------------------------------------------
\newcommand{\lstcode}{\ttfamily\color{qatextdark}}

\lstdefinelanguage{TypeScript}{
  keywords={typeof, new, true, false, catch, function, return, null, switch, var,
    if, in, while, do, else, case, break, const, let, interface, class, extends,
    implements, enum, export, import, this, super, constructor, public, private,
    protected, static, async, await, throw, throws, try, finally, void},
  ndkeywords={boolean, any, number, string, int, long, double},
  sensitive=true,
  comment=[l]{//},
  morecomment=[s]{/*}{*/},
  morestring=[b]',
  morestring=[b]"
}

\lstdefinestyle{r402code}{
  basicstyle=\lstcode\footnotesize,
  keywordstyle=\color{qacodekwl}\bfseries,
  ndkeywordstyle=\color{qaaccent2dark}\bfseries,
  stringstyle=\color{qacodestrl},
  commentstyle=\color{qacodecoml}\itshape,
  morecomment=[s][\color{qacodeannl}]{/**}{*/},
  backgroundcolor=\color{qapanellight},
  rulecolor=\color{qalinelight},
  frame=single,
  framesep=4pt,
  xleftmargin=6pt,
  xrightmargin=6pt,
  numbers=left,
  numberstyle=\tiny\color{qamuteddark},
  numbersep=8pt,
  keepspaces=true,
  columns=fullflexible,
  showstringspaces=false,
  showspaces=false,
  showtabs=false,
  breaklines=true,
  breakatwhitespace=false,
  captionpos=b,
  tabsize=2,
  escapeinside={(*@}{@*)},
  % accents UTF-8 dans le code (listings ne gere pas l'UTF-8 multi-octets)
  literate={é}{{\'e}}1 {è}{{\`e}}1 {ê}{{\^e}}1 {à}{{\`a}}1 {â}{{\^a}}1
           {ù}{{\`u}}1 {û}{{\^u}}1 {ô}{{\^o}}1 {î}{{\^i}}1 {ç}{{\c c}}1
           {É}{{\'E}}1 {À}{{\`A}}1 {€}{{\euro{}}}1 {…}{{\ldots}}1,
}
\lstset{style=r402code}

% ---------------------------------------------------------------
% BOITES : exercice, code, information
% ---------------------------------------------------------------
\tcbset{
  exobox/.style={colback=qapanellight, colframe=qaaccentdark, boxrule=0.6pt,
    arc=2pt, left=8pt, right=8pt, top=6pt, bottom=6pt,
    coltitle=white, colbacktitle=qaaccentdark, fonttitle=\bfseries,
    boxsep=2pt, breakable},
  infobox/.style={colback=white, colframe=qaaccent2dark, boxrule=0.6pt,
    leftrule=3pt, arc=0pt, left=8pt, right=8pt, top=5pt, bottom=5pt,
    coltitle=qaaccent2dark, colbacktitle=white, fonttitle=\bfseries,
    breakable},
  codebox/.style={colback=qatextdark, colframe=qalinelight, boxrule=0.4pt,
    arc=2pt, left=8pt, right=8pt, top=5pt, bottom=5pt, listing only,
    listing options={basicstyle=\ttfamily\small\color{white},
      breaklines=true, showstringspaces=false, columns=fullflexible,
      keepspaces=true, frame=none, numbers=none, backgroundcolor={}}, breakable}
}
\newtcolorbox[auto counter]{exobox}[2][]{exobox,
  title={Exercice~\thetcbcounter\ifstrempty{#2}{}{ --- #2}}, #1}
\newtcolorbox{infobox}[1][\`A retenir]{infobox, title={#1}}
\newtcblisting{codebox}{codebox}

% ---------------------------------------------------------------
% EN-TETE DE PAGE (fil d'Ariane + pagination)
% ---------------------------------------------------------------
\newcommand{\qaheadtitlevalue}{}
\newcommand{\qaheadtitle}[1]{\renewcommand{\qaheadtitlevalue}{#1}}

\pagestyle{fancy}
\fancyhf{}
\renewcommand{\headrulewidth}{0.4pt}
\fancyhead[L]{\small\color{qamuteddark}R4.02 --- \qaheadtitlevalue}
\fancyhead[R]{\small\color{qamuteddark}\thepage\,/\,\pageref{LastPage}}
\fancyfoot[C]{\footnotesize\color{qamuteddark}IUT de Vannes --- Universit\'e Bretagne Sud}
\fancypagestyle{plain}{\fancyhf{}\renewcommand{\headrulewidth}{0pt}%
  \fancyfoot[C]{\footnotesize\color{qamuteddark}IUT de Vannes --- Universit\'e Bretagne Sud}}

% ---------------------------------------------------------------
% BANDEAU DE TITRE (pas une page dediee, pour rester compact)
% \rtpheader{numero ou intitule court}{titre du TP}
% ---------------------------------------------------------------
\newcommand{\ryear}{2026/2027}
\newcommand{\rtpheader}[2]{%
  \thispagestyle{plain}%
  \begin{minipage}[t]{0.68\textwidth}
    {\small\color{qamuteddark}
      Universit\'e Bretagne Sud\\
      IUT de Vannes\\
      D\'epartement Informatique --- \ryear}
  \end{minipage}%
  \hfill
  \begin{minipage}[t]{0.28\textwidth}
    \raggedleft
    \includegraphics[width=2.6cm]{iut_logo.jpg}
  \end{minipage}

  \vspace{6mm}
  \begin{center}
    {\large\bfseries\color{qatextdark} BUT 2 --- R4.02. Qualit\'e de d\'eveloppement}\\[2mm]
    {\normalsize\color{qamuteddark} TP #1}\\[1mm]
    {\Large\bfseries\color{qaaccentdark} #2}\\[2mm]
    \raimention
  \end{center}
  \vspace{4mm}
  \qaheadtitle{#2}
}

\graphicspath{{../../latex/common/assets/}{./}}

% A ajuster avant diffusion
\newcommand{\rdatelimite}{la veille de la séance du TP4, à 23\,h\,59}

\begin{document}

\rtpheader{3}{Tests unitaires isolés avec Mockito}

\begin{infobox}[Objectifs]
\begin{itemize}[leftmargin=1.2em, itemsep=1pt]
  \item Tester un service sans base de données, sans réseau et sans dépendre de la date du jour.
  \item Utiliser Mockito : \texttt{@Mock}, \texttt{when ... thenReturn / thenThrow}, \texttt{verify}, \texttt{never}, \texttt{ArgumentCaptor}.
  \item Rendre testable une classe qui ne l'est pas, en injectant ses dépendances.
  \item Distinguer un test qui vérifie le comportement d'un test trop couplé à l'implémentation.
\end{itemize}
\end{infobox}

Continuez dans le dépôt de votre binôme, à partir de votre rendu du TP2. Mockito est déjà dans le \texttt{pom.xml} (\texttt{mockito-junit-jupiter}).

% ================================================================
\section{Le service à tester}

La classe \texttt{VirementService} (paquetage \texttt{fr.iutvannes.banque.service}) dépend de quatre collaborateurs, tous reçus par son constructeur (\texttt{CompteRepository} est dans le paquetage \texttt{persistance}) :

\begin{lstlisting}[language=Java]
public interface CompteRepository {
    Optional<Compte> findByNumero(String numero);
    Compte save(Compte compte);
    List<Compte> findAll();
    List<Compte> findByTitulaireOrderByNumero(String titulaire);
}
public interface TauxChangeClient {
    BigDecimal taux(Devise de, Devise vers);   // peut lever TauxIndisponibleException
}
public interface NotificationService {
    void notifier(String numeroCompte, String message);   // peut lever NotificationException
}

public VirementService(CompteRepository comptes, TauxChangeClient taux,
                       NotificationService notifications, Clock horloge)

public void virer(String numeroSource, String numeroCible, BigDecimal montant)
\end{lstlisting}

La méthode \texttt{virer} doit respecter les règles suivantes. Les règles du TP1 sur \texttt{Compte} (montant, découvert, plafond\ldots) s'appliquent aussi : c'est \texttt{Compte} qui les vérifie.

\begin{center}
\renewcommand{\arraystretch}{1.3}
\begin{tabular}{@{}>{\bfseries\color{qaaccentdark}}l >{\raggedright\arraybackslash}p{14cm}@{}}
\toprule
V1 & \textbf{Comptes.} Les deux comptes sont chargés par leur numéro. Si l'un d'eux n'existe pas : \texttt{CompteInconnuException}. Virer vers le même compte : \texttt{IllegalArgumentException}. \\
V2 & \textbf{Jour.} Aucun virement n'est accepté le dimanche (selon l'horloge, fuseau \texttt{Europe/Paris}) : \texttt{VirementImpossibleException}. \\
V3 & \textbf{Même devise.} La source est débitée du montant, la cible est créditée du même montant. Le service de change n'est \textbf{pas} appelé. \\
V4 & \textbf{Devises différentes.} La cible est créditée de \texttt{montant × taux(deviseSource, deviseCible)}, arrondi à 2 décimales (\texttt{RoundingMode.HALF\_EVEN}). Si le service de change est indisponible (\texttt{TauxIndisponibleException}) : \texttt{VirementImpossibleException}. \\
V5 & \textbf{Enregistrement.} Après un virement réussi, les deux comptes sont enregistrés. Un virement refusé (quelle qu'en soit la raison) ne modifie et n'enregistre \textbf{aucun} compte, et n'envoie aucune notification. \\
V6 & \textbf{Notification.} Si le montant est \textbf{strictement supérieur} à 1\,000,00\,\euro{}, le titulaire du compte source est notifié avec le message \texttt{"Virement de <montant> € effectué vers <numeroCible>"}. Si l'envoi de la notification échoue (\texttt{NotificationException}), le virement \textbf{reste valide} : l'exception n'est pas propagée. \\
\bottomrule
\end{tabular}
\end{center}

% ================================================================
\section{Tester \texttt{VirementService} avec Mockito (40 min)}

\begin{exobox}{Mise en place}
Créez \texttt{src/test/java/fr/iutvannes/banque/service/VirementServiceTest.java} :
\begin{itemize}[leftmargin=1.2em]
  \item annotez la classe avec \texttt{@ExtendWith(MockitoExtension.class)} ;
  \item déclarez des \texttt{@Mock} pour le repository, le service de change et le service de notification ;
  \item créez une horloge fixe un jour de semaine, par exemple :
\begin{codebox}
Clock.fixed(Instant.parse("2027-03-01T10:00:00Z"),
            ZoneId.of("Europe/Paris"))
\end{codebox}
  \item construisez le service dans un \texttt{@BeforeEach}.
\end{itemize}
Écrivez une petite méthode utilitaire qui crée un compte et programme \texttt{findByNumero} pour le renvoyer : vous allez en avoir besoin dans presque tous les tests.
\end{exobox}

\begin{exobox}{Une règle, un groupe de tests}
Écrivez les tests des règles V1 à V6, regroupés par règle avec \texttt{@Nested}. Pour chaque règle, testez au moins un cas nominal, un cas limite et un cas d'erreur. Quelques indications :
\begin{enumerate}[leftmargin=1.4em]
  \item \textbf{V1} : un repository qui renvoie \texttt{Optional.empty()}. Pensez à V5 : rien ne doit être enregistré (\texttt{verify(..., never())}).
  \item \textbf{V2} : il faut une \textbf{autre} horloge, fixée un dimanche. Où se trouve exactement la limite entre samedi et dimanche ? Attention au fuseau horaire : \texttt{2027-03-06T23:30:00Z}, c'est quel jour à Paris ?
  \item \textbf{V3} : vérifiez les soldes (état) \textbf{et} que le service de change n'a pas été appelé (\texttt{verifyNoInteractions}).
  \item \textbf{V4} : programmez le taux avec \texttt{when(...).thenReturn(...)}. Choisissez un taux et un montant qui obligent à arrondir. Pour la panne, utilisez \texttt{thenThrow}.
  \item \textbf{V5} : après un virement réussi, vérifiez que chaque compte est enregistré. Après un solde insuffisant (levé par \texttt{Compte}), vérifiez que rien n'est enregistré ni notifié.
  \item \textbf{V6} : testez les montants 1\,000,00, 1\,000,01 et un montant plus élevé. Récupérez le message avec un \texttt{ArgumentCaptor} et vérifiez son contenu exact.
  \item \textbf{V6, échec de notification} : \texttt{notifier} renvoie \texttt{void}, on ne peut donc pas écrire \texttt{when(...)}. Utilisez la forme suivante ; le virement doit réussir quand même.
\begin{codebox}
doThrow(new NotificationException("SMS indisponible"))
    .when(notifications).notifier(any(), any());
\end{codebox}
\end{enumerate}
\end{exobox}

\begin{infobox}[Si Mockito proteste]
\begin{itemize}[leftmargin=1.2em, itemsep=1pt]
  \item \texttt{UnnecessaryStubbingException} : vous avez programmé une réponse que le test n'utilise pas. Supprimez-la : elle rend le test trompeur.
  \item \texttt{InvalidUseOfMatchersException} : vous mélangez des matchers (\texttt{any()}) et des valeurs brutes dans un même appel. Entourez les valeurs avec \texttt{eq(...)}.
  \item \texttt{WantedButNotInvoked} / \texttt{TooManyActualInvocations} : le message indique les appels réellement reçus par le mock ; lisez-le, il contient souvent la réponse.
\end{itemize}
\end{infobox}

\begin{exobox}{Corriger les bugs}
Comme au TP1, le code fourni contient des bugs. Pour chacun : notez dans \texttt{docs/tp3.md} la règle violée et le test qui le révèle, corrigez le code, relancez \textbf{toute} la suite, puis faites un commit.

Question : les tests Cucumber du TP2 utilisent aussi \texttt{VirementService}. Pourquoi n'ont-ils pas détecté ces bugs ? Qu'en concluez-vous sur la complémentarité des deux niveaux de tests ?
\end{exobox}

% ================================================================
\section{Rendre une classe testable (30 min)}

La classe \texttt{AlerteDecouvertJob} est exécutée chaque matin. Elle prévient par SMS les titulaires des comptes à découvert, sauf le week-end (les alertes du samedi et du dimanche sont envoyées le lundi). Elle a été écrite rapidement et n'a aucun test :

\begin{lstlisting}[language=Java]
public class AlerteDecouvertJob {

    public int executer() {
        CompteRepository comptes = new JdbcCompteRepository(
            "jdbc:postgresql://localhost:5432/banque", "banque", "secret");
        NotificationService sms =
            new SmsNotificationService("https://api.sms.example", "CLE-API");

        DayOfWeek jour = LocalDate.now().getDayOfWeek();
        if (jour == DayOfWeek.SATURDAY || jour == DayOfWeek.SUNDAY) {
            return 0;
        }
        int nbAlertes = 0;
        for (Compte compte : comptes.findAll()) {
            if (compte.getSolde().signum() < 0) {
                sms.notifier(compte.getNumero(),
                             "Votre compte est à découvert le " + LocalDate.now());
                nbAlertes++;
            }
        }
        return nbAlertes;
    }
}
\end{lstlisting}

\begin{exobox}{Refactoriser, puis tester}
\begin{enumerate}[leftmargin=1.4em]
  \item Dans \texttt{docs/tp3.md}, listez tout ce qui empêche d'écrire un test unitaire de cette classe. Pour chaque point, dites ce qui se passerait si on lançait quand même un test.
  \item Refactorisez la classe pour que ses dépendances (repository, notifications, horloge) soient \textbf{injectées par le constructeur}. Le comportement ne doit pas changer. Où les vraies implémentations seront-elles créées désormais ?
  \item Écrivez ses tests avec Mockito : aucun compte, aucun compte à découvert, deux comptes à découvert sur trois, compte à solde exactement nul, samedi, dimanche, lundi. Vérifiez la valeur renvoyée \textbf{et} les notifications envoyées (destinataire et texte, date comprise).
  \item Un SMS échoue (\texttt{NotificationException}) pour le premier compte : que fait la classe actuellement ? Est-ce souhaitable ? Écrivez le test qui fixe le comportement que vous jugez correct (les autres titulaires doivent-ils être prévenus ?), puis adaptez le code.
\end{enumerate}
\end{exobox}

% ================================================================
\section{Des tests trop couplés ? (20 min)}

\begin{exobox}{Un refactoring qui casse des tests}
Un collègue propose une petite optimisation dans \texttt{VirementService} : remplacer les deux appels à \texttt{comptes.save(...)} par un seul appel à une nouvelle méthode \texttt{comptes.saveAll(List.of(source, cible))}. Ajoutez-la à l'interface \texttt{CompteRepository} (\texttt{List<Compte> saveAll(Iterable<Compte> comptes);} : Spring Data la fournit) et à \texttt{CompteRepositoryEnMemoire} (paquetage de test \texttt{doublures}).
\begin{enumerate}[leftmargin=1.4em]
  \item Faites la modification et relancez vos tests. Lesquels échouent ?
  \item Le comportement observable du service a-t-il changé pour un utilisateur de la banque ? Les tests qui échouent signalaient-ils un vrai problème ?
  \item Réécrivez ces tests pour qu'ils vérifient ce qui compte (les comptes enregistrés et leurs soldes), sans dépendre de la manière exacte de les enregistrer. Une piste : remplacer le mock du repository par un \textbf{fake} en mémoire (\texttt{CompteRepositoryEnMemoire}, fourni dans le paquetage de test \texttt{doublures}) dans certains tests.
  \item En deux ou trois phrases dans \texttt{docs/tp3.md} : dans quels cas un \texttt{verify} est-il indispensable, et dans quels cas vaut-il mieux vérifier un état ?
\end{enumerate}
\end{exobox}

% ================================================================
\clearpage
\section{Rendu}

\begin{infobox}[À rendre avant \rdatelimite]
Sur le dépôt du binôme, un \textbf{tag} \texttt{tp3} pointant sur le commit rendu :
\begin{codebox}
git tag tp3
git push origin tp3
\end{codebox}
Le commit doit contenir :
\begin{itemize}[leftmargin=1.2em, itemsep=1pt]
  \item \texttt{VirementServiceTest.java} (règles V1 à V6) et \texttt{VirementService.java} corrigé ;
  \item \texttt{AlerteDecouvertJob.java} refactorisée et \texttt{AlerteDecouvertJobTest.java} ;
  \item les tests réécrits après le passage à \texttt{saveAll} ;
  \item \texttt{docs/tp3.md} : journal des bugs, complémentarité avec Cucumber, analyse de \texttt{AlerteDecouvertJob}, réflexion sur \texttt{verify} ;
  \item une suite complète qui passe : \texttt{./mvnw test}.
\end{itemize}
\end{infobox}

\end{document}
